\documentclass[11pt,a4paper]{article}

\usepackage[utf8]{inputenc}
\usepackage[T1]{fontenc}
\usepackage{amsmath}
\usepackage{amssymb}
\usepackage{booktabs}
\usepackage{hyperref}
\usepackage{graphicx}
\usepackage{geometry}
\usepackage{xcolor}
\usepackage{listings}
\usepackage{array}
\usepackage{caption}
\usepackage{tikz}
\usepackage{tcolorbox}

\geometry{margin=1in}

\definecolor{codeblue}{rgb}{0.1,0.2,0.6}
\definecolor{codegray}{rgb}{0.5,0.5,0.5}
\definecolor{codepurple}{rgb}{0.58,0,0.82}
\definecolor{backcolour}{rgb}{0.97,0.97,0.97}
\definecolor{alertred}{rgb}{0.8,0.0,0.0}
\definecolor{okgreen}{rgb}{0.0,0.5,0.0}

\lstset{
  language=Python,
  backgroundcolor=\color{backcolour},
  basicstyle=\ttfamily\small,
  keywordstyle=\color{codeblue}\bfseries,
  commentstyle=\color{codegray}\itshape,
  stringstyle=\color{codepurple},
  breaklines=true,
  frame=single,
  numbers=left,
  numberstyle=\tiny\color{codegray},
  showstringspaces=false,
}

\tcbuselibrary{skins,breakable}
\newtcolorbox{insight}{
  colback=yellow!10!white,
  colframe=orange!70!black,
  title=Key Insight,
  fonttitle=\bfseries,
}
\newtcolorbox{bugbox}{
  colback=red!5!white,
  colframe=alertred,
  title=Code Bug,
  fonttitle=\bfseries,
}

\title{AdamW Optimizer Memory Overhead: \\
       The 2X vs.\ 3X Controversy \\[0.5em]
       \large Focusing on Peak (Maximum) Memory Usage}
\author{}
\date{May 2026}

\begin{document}
\maketitle

\begin{abstract}
A recurring question in GPU memory analysis is whether the AdamW optimizer
occupies \emph{two} extra copies of the model parameters (2X) or
\emph{three} extra copies (3X) at peak.  The 2X figure counts the two
persistent state tensors---\texttt{exp\_avg} and \texttt{exp\_avg\_sq}.
The 3X figure arises because \texttt{optimizer.step()} internally computes
$\sqrt{v_t}$ via a \emph{temporary} tensor that is live \emph{during} the
step but freed immediately after.  Since our analysis focuses on
\textbf{maximum (peak) memory usage}, and the temporary tensor does occupy
GPU memory at peak, the correct number for peak analysis is \textbf{3X}.
This document traces the evidence through the repository code and explains
the peak-detection algorithm.
\end{abstract}

%=============================================================================
\tableofcontents
\newpage

%=============================================================================
\section{Location of AdamW in the Codebase}
\label{sec:location}
%=============================================================================

The \texttt{torch.optim.AdamW} optimizer appears in two scripts:

\begin{table}[h]
\centering
\caption{Scripts that instantiate AdamW.}
\label{tab:scripts}
\begin{tabular}{lll}
\toprule
\textbf{File} & \textbf{Purpose} & \textbf{Relevant Line} \\
\midrule
\texttt{why\_adam\_times\_3\_here/iiiandadam.py}
  & Isolated micro-benchmark; proves the 3X peak
  & \texttt{optimizer = torch.optim.AdamW(\ldots)} \\
\texttt{gpt2\_gpipe\_memory\_profile.py}
  & Full GPT-2 pipeline-parallel training profiler
  & \texttt{optimizer = torch.optim.AdamW(\ldots)} \\
\bottomrule
\end{tabular}
\end{table}

Both scripts use the same API (\texttt{torch.optim.AdamW}), so the memory
behaviour is identical in both.

%=============================================================================
\section{Background: AdamW Update Formula}
\label{sec:formula}
%=============================================================================

Given a parameter tensor $\theta \in \mathbb{R}^n$ with gradient $g_t$ at
step $t$, AdamW performs:
\begin{align}
  m_t &= \beta_1 m_{t-1} + (1 - \beta_1)\, g_t
        \quad\text{(first moment, \texttt{exp\_avg})} \label{eq:m} \\
  v_t &= \beta_2 v_{t-1} + (1 - \beta_2)\, g_t^2
        \quad\text{(second moment, \texttt{exp\_avg\_sq})} \label{eq:v} \\
  \hat{m}_t &= m_t / (1 - \beta_1^t)
        \quad\text{(bias-corrected first moment)} \label{eq:mhat} \\
  \hat{v}_t &= v_t / (1 - \beta_2^t)
        \quad\text{(bias-corrected second moment)} \label{eq:vhat} \\
  \theta_t &\leftarrow \theta_{t-1}
            - \alpha \left(\frac{\hat{m}_t}{\sqrt{\hat{v}_t} + \epsilon}
            + \lambda\,\theta_{t-1}\right) \label{eq:update}
\end{align}

The bias-correction for the second moment is absorbed into the learning
rate in PyTorch's implementation, but the critical operation is
\begin{equation}
  \mathtt{denom} = \sqrt{v_t} + \epsilon
  \label{eq:denom}
\end{equation}
which is computed as a \textbf{new tensor} before $\theta$ is updated.

%=============================================================================
\section{Memory Terminology: 1X, 2X, 3X}
\label{sec:terminology}
%=============================================================================

Throughout this document we use the following shorthand.

\begin{table}[h]
\centering
\caption{Memory multiplier notation. $P$ denotes the total size (in bytes)
of \emph{all model parameters} on one GPU, using FP32.}
\label{tab:notation}
\begin{tabular}{cll}
\toprule
\textbf{Notation} & \textbf{Meaning} & \textbf{Example (FP32)} \\
\midrule
$1\mathrm{X} = P$   & One full copy of model parameters & \\
$k\mathrm{X}$        & $k$ copies of model parameter size
                     & $2\mathrm{X} = 2P$ bytes \\
\bottomrule
\end{tabular}
\end{table}

\noindent\textbf{Important:} ``extra beyond model parameters'' means we
count $m_t$, $v_t$, $\sqrt{v_t}$, and gradients as multiples of $P$, but
\emph{not} $P$ itself.

%=============================================================================
\section{The 2X Claim: Two Persistent Optimizer State Tensors}
\label{sec:2x}
%=============================================================================

When \texttt{optimizer.step()} is called for the \emph{first} time, AdamW
lazily allocates the state tensors:

\begin{itemize}
  \item \textbf{\texttt{exp\_avg}} ($m_t$): same shape and dtype as the
    parameter $\theta$. Size $= |\theta|\cdot\beta = P$ per parameter.
  \item \textbf{\texttt{exp\_avg\_sq}} ($v_t$): same shape and dtype.
    Size $= P$ per parameter.
\end{itemize}

These two tensors persist for the entire training run.  Their combined size
is:
\begin{equation}
  \text{Persistent AdamW states} = |\mathtt{exp\_avg}| + |\mathtt{exp\_avg\_sq}|
  = P + P = 2P = 2\mathrm{X}
  \label{eq:2x}
\end{equation}

\medskip
\noindent Evidence from \texttt{iiiandadam.py} (see Section~\ref{sec:iiiandadam}):
\begin{quote}
\textit{``記憶體會暴增！因為 Adam 分配了 \texttt{exp\_avg} 和
\texttt{exp\_avg\_sq} (2倍的模型大小)''}
\end{quote}
The script observes this with \texttt{torch.cuda.memory\_allocated()} after
\texttt{step()} completes.  Because \texttt{memory\_allocated()} reports
the \emph{currently allocated} (not the \emph{peak during step}) bytes, it
correctly shows $+2P$ after step.

%=============================================================================
\section{The 3X Peak: The Temporary \texorpdfstring{$\sqrt{v_t}$}{sqrt(v)}
         Buffer}
\label{sec:3x}
%=============================================================================

\subsection{Where the Extra Allocation Comes From}

PyTorch's AdamW implementation (both the Python loop path and the C++
kernel) computes the denominator as:
\begin{lstlisting}[language=Python, caption={Simplified AdamW step kernel (illustrative).}]
# Inside torch.optim.AdamW.step()  (simplified)
exp_avg.mul_(beta1).add_(grad, alpha=1 - beta1)        # in-place update of m
exp_avg_sq.mul_(beta2).addcmul_(grad, grad, value=1-beta2)  # in-place update of v

# KEY LINE: creates a NEW temporary tensor for sqrt(v)
denom = (exp_avg_sq.sqrt() / math.sqrt(bias_correction2)).add_(eps)

step_size = lr / bias_correction1
param.addcdiv_(exp_avg, denom, value=-step_size)       # update theta in-place
# denom goes out of scope and is freed here
\end{lstlisting}

The call \texttt{exp\_avg\_sq.sqrt()} returns a \textbf{new tensor}
(it does \emph{not} modify \texttt{exp\_avg\_sq} in-place).  This new
tensor --- let us call it \texttt{sqrt\_buf} --- has the same shape and
dtype as $v_t$ and thus occupies $P$ bytes.

\subsection{Memory Picture at Peak vs.\ After Step}

\begin{table}[h]
\centering
\caption{Live tensors at two critical moments during \texttt{optimizer.step()}.
$G = P$ because gradients have the same shape as parameters.}
\label{tab:peak_vs_after}
\begin{tabular}{llcc}
\toprule
\textbf{Tensor} & \textbf{Description}
  & \textbf{Live during step} & \textbf{Live after step} \\
\midrule
$\theta$              & Model parameters          & \checkmark & \checkmark \\
$g_t$                 & Gradients                 & \checkmark & \checkmark \\
\texttt{exp\_avg}     & First moment $m_t$        & \checkmark & \checkmark \\
\texttt{exp\_avg\_sq} & Second moment $v_t$       & \checkmark & \checkmark \\
\texttt{sqrt\_buf}    & $\sqrt{v_t}$ (temporary)  & \checkmark & \\
\midrule
\textbf{Extra beyond $\theta$} (in units of $P$)
  & & $\mathbf{4P}$
  & $\mathbf{3P}$ \\
\textbf{Adam-specific extra} (excluding grads)
  & & $\mathbf{3P = 3\mathrm{X}}$
  & $\mathbf{2P = 2\mathrm{X}}$ \\
\bottomrule
\end{tabular}
\end{table}

\begin{insight}
Because our analysis focuses on \textbf{maximum (peak) memory}, the
correct multiplier for AdamW's overhead is \textbf{3X}, not 2X.  The
2X figure is the \emph{sustained} post-step cost; the 3X figure is the
\emph{transient peak} cost that occurs while \texttt{step()} is executing.
\end{insight}

%=============================================================================
\section{Code Analysis}
\label{sec:code}
%=============================================================================

\subsection{Isolated Micro-Benchmark: \texttt{iiiandadam.py}}
\label{sec:iiiandadam}

The file \texttt{why\_adam\_times\_3\_here/iiiandadam.py} is a self-contained
script designed to isolate and observe AdamW's memory behaviour step by
step.

\begin{lstlisting}[caption={\texttt{why\_adam\_times\_3\_here/iiiandadam.py} (full script).}]
import torch
import torch.nn as nn

torch.cuda.memory._record_memory_history()   # start tracing ALL alloc/free events

device = torch.device("cuda")

# Step 1: load model  ->  allocates P bytes
model = nn.Linear(10000, 1000).to(device)
# weight: 10000*1000*4 = 40,000,000 B  (~38.15 MB)
# bias:      1000*4   =      4,000 B  (  ~0.004 MB)
# Total P ~ 38.15 MB

# Step 2: declare optimizer  ->  no GPU memory yet (lazy state init)
optimizer = torch.optim.AdamW(model.parameters(), lr=1e-3)

# Step 3: forward + backward  ->  gradients allocated (another P)
dummy_input = torch.randn(16, 10000).to(device)
output = model(dummy_input)
loss = output.sum()
loss.backward()
# Alive: theta (P) + grads (P) = 2P

# Step 4: first optimizer.step()
#   SUSTAINED after: theta (P) + grads (P) + exp_avg (P) + exp_avg_sq (P) = 4P
#   PEAK DURING:     theta (P) + grads (P) + exp_avg (P) + exp_avg_sq (P)
#                              + sqrt_buf (P) = 5P
optimizer.step()
# torch.cuda.memory_allocated() reports 4P here (sustained, sqrt_buf freed)

snapshot_file = "gpu_memory_snapshot-adam.pickle"
torch.cuda.memory._dump_snapshot(snapshot_file)   # captures full event trace
torch.cuda.memory._record_memory_history(enabled=None)
\end{lstlisting}

\paragraph{What \texttt{memory\_allocated()} shows vs.\ what the snapshot reveals.}
The \texttt{print} statement after \texttt{optimizer.step()} calls
\texttt{torch.cuda.memory\_allocated()}, which reports the bytes
\emph{currently allocated} at that moment.  By then, \texttt{sqrt\_buf}
has already been freed, so the readout is $4P$ (not $5P$).

To observe the true peak --- which includes \texttt{sqrt\_buf} --- one
must replay the full event trace stored in the snapshot using the
\texttt{ref\_find\_max/MemoryViz.js} algorithm
(Section~\ref{sec:peak_detection}).

\subsection{GPT-2 Pipeline Profiler: \texttt{gpt2\_gpipe\_memory\_profile.py}}
\label{sec:gpt2}

The full training script follows the same pattern:

\begin{lstlisting}[caption={AdamW usage in \texttt{gpt2\_gpipe\_memory\_profile.py} (abridged).},firstnumber=296]
# Optimizer -- same API as iiiandadam.py
optimizer = torch.optim.AdamW(pipe_model.parameters(), lr=LEARNING_RATE)

# Training loop
for step in range(NUM_STEPS):
    logits = pipe_model(input_ids)           # forward through pipeline
    loss   = nn.functional.cross_entropy(...)

    # Backward group
    optimizer.zero_grad()                    # clear stale gradients
    loss.backward()                          # create fresh gradients
    optimizer.step()                         # <- peak here (includes sqrt_buf)
\end{lstlisting}

\paragraph{Why \texttt{zero\_grad()} does not affect the peak.}
\texttt{optimizer.zero\_grad()} is called \emph{before} \texttt{loss.backward()},
so it merely clears \emph{stale} gradients from the previous iteration.
Fresh gradients of size $P$ are re-created by \texttt{loss.backward()}.
The peak still occurs during \texttt{optimizer.step()} with fresh gradients
live:
\begin{equation}
  M_{\text{peak at step}} = P_\theta + P_g + P_{m} + P_{v} + P_{\sqrt{v}}
  = 5P
  \label{eq:5p}
\end{equation}
where ``extra due to Adam'' $= P_m + P_v + P_{\sqrt{v}} = 3P = 3\mathrm{X}$.

\paragraph{Memory accounting in the full GPT-2 run (GPU 0).}
From the companion document \texttt{memory\_analysis.tex}, the 96
optimizer-state blocks (category B) sum to $226{,}811{,}904$ bytes $= 2
\times 113{,}405{,}952$ bytes $= 2 \times P_{\text{GPU0}}$.  This is the
\emph{sustained} 2X cost after step completes.  The 3X peak cost includes
the transient \texttt{sqrt\_buf} which is not present in the final snapshot
(it is freed within \texttt{step()}).

%=============================================================================
\section{Peak Memory Detection Methodology}
\label{sec:peak_detection}
%=============================================================================

\subsection{Recording and Exporting the Trace}

\begin{enumerate}
  \item \textbf{Start recording.} \\
        \texttt{torch.cuda.memory.\_record\_memory\_history()} registers a
        callback that fires on every \texttt{cudaMalloc} and \texttt{cudaFree}
        call.  Every event (address, size, stack trace, timestamp) is stored
        in a ring buffer.

  \item \textbf{Export the trace.} \\
        \texttt{torch.cuda.memory.\_dump\_snapshot(path)} serialises the
        ring buffer to a \texttt{.pickle} file.  The exported data structure
        contains a \texttt{device\_traces} list --- one ordered list of
        allocation events per GPU.

  \item \textbf{Replay and find peak.} \\
        \texttt{ref\_find\_max/MemoryViz.js} (function
        \texttt{process\_alloc\_data}) replays the event trace to
        reconstruct the memory-over-time curve and locate the exact peak.
\end{enumerate}

\subsection{\texttt{MemoryViz.js} Peak-Finding Algorithm}
\label{sec:memviz}

The algorithm (from \texttt{ref\_find\_max/MemoryViz.js}) works as follows:

\begin{lstlisting}[language=JavaScript, caption={Simplified peak-finding loop
  from \texttt{MemoryViz.js}.}]
let total_mem = 0;
let max_size  = 0;
let peak_elem = null;
const max_at_time = [];

for (const elem of actions) {              // replay alloc/free events in order
  const size = elements[elem].size;
  const idx  = current.findLastIndex(x => x === elem);

  if (idx === -1) {
    // first occurrence -> allocation
    add_allocation(elem);                  // total_mem += size
    advance(1);                            // max_at_time.push(total_mem)
  } else {
    // second occurrence -> free
    advance(1);
    remove_allocation(elem);              // total_mem -= size
  }

  if (total_mem > max_size) {
    max_size  = total_mem;
    peak_elem = elem;                      // record which event caused the peak
  }
}

// Identify peak timestep
let peak_timestep = max_at_time.indexOf(Math.max(...max_at_time));

// Collect all blocks alive at peak_timestep
const peak_allocs = data.filter(d =>
  d.timesteps[0] <= peak_timestep &&
  peak_timestep  <= d.timesteps[d.timesteps.length - 1]
);
\end{lstlisting}

\paragraph{Key property.}
Because the trace records every \texttt{alloc} and \texttt{free}---including
the transient \texttt{sqrt\_buf} allocation that appears and disappears
\emph{inside} \texttt{optimizer.step()}---the algorithm correctly captures
that the peak occurred while \texttt{sqrt\_buf} was live, even though
\texttt{sqrt\_buf} does not appear in the final snapshot.

The output (\texttt{peak\_alloc\_events\_GPU0.json}) lists every tensor
alive at the peak timestep.  For the full GPT-2 run this comprises
\textbf{282 blocks totalling 625.77\,MB} on GPU\,0.

%=============================================================================
\section{Code Bug: Unnecessary Per-GPU Loop in Memory Recording}
\label{sec:bug}
%=============================================================================

\begin{bugbox}
\texttt{gpt2\_gpipe\_memory\_profile.py} contains two helper functions that
iterate over GPUs before calling the memory-recording API.  This loop is
\textbf{unnecessary} in modern PyTorch and should be removed.
\end{bugbox}

\subsection{Current (Buggy) Code}

\begin{lstlisting}[caption={Current implementation in
  \texttt{gpt2\_gpipe\_memory\_profile.py} (lines 248--266).}]
def start_memory_history_all_gpus():
    """Start recording memory history on ALL GPUs."""
    print("[INFO] Starting memory history recording on all GPUs...")
    for i in range(NUM_GPUS):              # <-- loop is unnecessary
        with torch.cuda.device(i):         # <-- context has no effect on API
            torch.cuda.memory._record_memory_history(
                max_entries=1048576
            )

def stop_memory_history_all_gpus():
    """Stop recording memory history on ALL GPUs."""
    for i in range(NUM_GPUS):              # <-- loop is unnecessary
        with torch.cuda.device(i):
            torch.cuda.memory._record_memory_history(enabled=None)
    print("[INFO] Stopped memory history recording on all GPUs.")
\end{lstlisting}

\subsection{Why the Loop is Wrong}

\texttt{torch.cuda.memory.\_record\_memory\_history()} is a \textbf{global
process-level} call in modern PyTorch.  It enables memory event recording
for \emph{all} CUDA devices at once.  The \texttt{with torch.cuda.device(i)}
context manager sets the default device for operations such as tensor
creation, but it does \textbf{not} scope the memory recording API to a
single GPU.

Consequently, the loop calls the same global function $\mathtt{NUM\_GPUS}$
times.  Each call overwrites the previous setting, so only the behaviour
of the \emph{last} iteration actually takes effect.  This is not harmful
in practice (the last call enables recording on all GPUs as desired), but
it is misleading and wasteful.

\subsection{Corrected Code}

\begin{lstlisting}[caption={Corrected implementation.}]
def start_memory_history_all_gpus():
    """Start recording memory history (global, covers all GPUs)."""
    print("[INFO] Starting memory history recording on all GPUs...")
    torch.cuda.memory._record_memory_history(max_entries=1048576)

def stop_memory_history_all_gpus():
    """Stop recording memory history."""
    torch.cuda.memory._record_memory_history(enabled=None)
    print("[INFO] Stopped memory history recording on all GPUs.")
\end{lstlisting}

Note that \texttt{dump\_memory\_snapshot\_all\_gpus()} is \emph{correctly}
implemented: it calls \texttt{torch.cuda.memory.\_snapshot()} \emph{once}
to obtain the global snapshot and then filters it per GPU.

%=============================================================================
\section{Summary: 2X vs.\ 3X}
\label{sec:summary}
%=============================================================================

\begin{table}[h]
\centering
\caption{AdamW memory overhead summary.  ``Extra'' is relative to model
parameters $P$.}
\label{tab:summary}
\begin{tabular}{lccl}
\toprule
\textbf{Scenario} & \textbf{Extra (AdamW only)} & \textbf{Total on GPU} & \textbf{When} \\
\midrule
Declared, not yet stepped & $0$   & $P$          & after \texttt{AdamW(\ldots)} \\
After \texttt{backward()}  & $0$   & $2P$         & grads live, no states yet \\
\textbf{Sustained post-step} & $\mathbf{2P}$ & $4P$ & \texttt{exp\_avg} + \texttt{exp\_avg\_sq} \\
\textbf{Peak during step}    & $\mathbf{3P}$ & $5P$ & $+$ transient \texttt{sqrt\_buf} \\
\midrule
\multicolumn{4}{l}{\textbf{Correct answer for peak analysis: 3X}} \\
\bottomrule
\end{tabular}
\end{table}

\begin{equation}
\boxed{
  M_{\text{peak}}^{\text{AdamW}} = P + \underbrace{P}_{\text{grads}}
  + \underbrace{P + P}_{\text{exp\_avg} + \text{exp\_avg\_sq}}
  + \underbrace{P}_{\sqrt{v_t}\text{ (temp)}}
  = 5P
}
\label{eq:peak_total}
\end{equation}

\noindent where the \emph{AdamW-specific} extra (beyond $\theta$ and $g$) is
$3P = 3\mathrm{X}$.

\medskip
\noindent\textbf{Why \texttt{zero\_grad()} is irrelevant to the peak.}
\texttt{optimizer.zero\_grad()} cannot reduce $M_\text{peak}$ because the
peak happens \emph{inside} \texttt{optimizer.step()}, by which time
\texttt{loss.backward()} has already created fresh gradients.  Whether
\texttt{zero\_grad()} was called before or after the backward pass does not
change what is live at the step-internal peak.

%=============================================================================
\section{Connection to the Full GPT-2 Analysis}
\label{sec:connection}
%=============================================================================

The companion document \texttt{memory\_analysis.tex} analyses the peak
GPU\,0 memory of the full GPT-2 GPipe training run.  The 282 blocks are
split into:

\begin{itemize}
  \item \textbf{(A) Parameters}: $4 \times P_\text{block} = 113{,}405{,}952$\,B
  \item \textbf{(B) AdamW states}: $2 \times 4 \times P_\text{block} = 226{,}811{,}904$\,B
        (the \emph{sustained} 2X cost)
  \item \textbf{(F) Gradients}: $4 \times P_\text{block} = 113{,}405{,}952$\,B
\end{itemize}

The transient \texttt{sqrt\_buf} (3rd extra $P$) is \emph{not} present in
the final snapshot used for that analysis, because the snapshot was taken
\emph{after} the training loop completed.  However, during any training
step, the true peak was $5P$ per partition, consistent with the 3X
conclusion above.

%=============================================================================
\section{Why the Reported 282-Block Peak Is Not the True Maximum}
\label{sec:notmax}
%=============================================================================

The companion document \texttt{memory\_analysis.tex} identifies \textbf{282 live
blocks totalling 625.77\,MB} on GPU\,0 as the peak memory snapshot.
This section presents concrete numerical evidence---drawn directly from
\texttt{memory\_reports\_org\_backup/} and
\texttt{ref\_find\_max/MemoryViz.js}---showing that $625.77\,\mathrm{MB}$ is
\textbf{not} the true maximum, and explains the two logical errors responsible
for the underestimate.

%─────────────────────────────────────────────────────────────────────────────
\subsection{Empirical Rebuttal: the Final Snapshot Alone Exceeds 625.77\,MB}
\label{sec:simplebound}
%─────────────────────────────────────────────────────────────────────────────

\begin{insight}
If $625.77\,\mathrm{MB}$ were the true peak, then every memory state
\emph{after} that moment must satisfy $M \leq 625.77\,\mathrm{MB}$.
But replaying \texttt{device\_traces[0]} in
\texttt{memory\_reports\_org\_backup/memory\_snapshot\_gpu0\_final.pickle}
to its last event (the moment the snapshot was taken, i.e.\ the
\emph{post-training final state}) yields:
\begin{equation}
  M_{\text{final}} = 987{,}414{,}528\,\text{B} = 941.67\,\text{MB},
  \quad 155\text{ blocks}
  \label{eq:finalstate}
\end{equation}
Since $941.67\,\mathrm{MB} > 625.77\,\mathrm{MB}$, the reported ``peak''
is falsified by the final state alone.
\end{insight}

\noindent The true peak, found by locating the global maximum of the
cumulative-allocation curve over all 40{,}156 events in
\texttt{device\_traces[0]}, is:
\begin{equation}
  \boxed{M_{\text{true peak}} = 1{,}230{,}004{,}224\,\text{B}
         = 1173.02\,\text{MB} \quad \text{at event index 8015}}
  \label{eq:truepeak}
\end{equation}
at which moment \textbf{193 blocks} are simultaneously alive.

\begin{table}[h]
\centering
\caption{Three memory figures derived from the same pickle file.}
\label{tab:three_figures}
\begin{tabular}{llrl}
\toprule
\textbf{Figure} & \textbf{Source} & \textbf{Value} & \textbf{Status} \\
\midrule
Reported ``peak'' from JSON   & \texttt{peak\_alloc\_events\_GPU0.json}
  & 625.77\,MB (282 blocks) & \textcolor{alertred}{\textbf{Wrong}} \\
Post-training final state     & trace replay, last event
  & 941.67\,MB (155 blocks) & Lower bound on true peak \\
True peak                     & trace replay, global max
  & \textbf{1173.02\,MB} (193 blocks) & \textcolor{okgreen}{\textbf{Correct}} \\
\bottomrule
\end{tabular}
\end{table}

%─────────────────────────────────────────────────────────────────────────────
\subsection{Logical Error 1: Snapshot Taken After the Training Loop}
\label{sec:err1timing}
%─────────────────────────────────────────────────────────────────────────────

In \texttt{gpt2\_gpipe\_memory\_profile.py} the snapshot is dumped
\emph{after} the training loop completes:

\begin{lstlisting}[caption={Incorrect snapshot placement in
    \texttt{gpt2\_gpipe\_memory\_profile.py} (abridged).}]
for step in range(NUM_STEPS):          # 10 training steps
    logits  = pipe_model(input_ids)    # forward
    loss    = cross_entropy(...)
    optimizer.zero_grad()
    loss.backward()                    # backward + recompute
    optimizer.step()                   # <-- TRUE PEAK IS HERE (inside each step)

# !! snapshot taken AFTER the loop -- all transient tensors already freed !!
snapshot_paths = dump_memory_snapshot_all_gpus("final")
\end{lstlisting}

When the snapshot is captured, every transient forward/backward tensor
(GPipe activation recomputations, intermediate attention buffers, the AdamW
\texttt{sqrt\_buf}) has already been freed.
The \emph{current-state} view in \texttt{snapshot.segments} therefore
reflects only the \emph{persistent} post-training memory
($\theta + g + m_t + v_t = 4P$), not the intra-step peak.

\medskip
\noindent\textbf{Why device\_traces doesn't save the methodology:}
Although \texttt{\_record\_memory\_history} captures all 40{,}156 alloc/free
events in \texttt{device\_traces[0]} (well within the 1\,M-entry ring
buffer\footnote{The ring buffer is $1{,}048{,}576$ entries; 40{,}156 GPU-0
events means \textbf{no overflow occurred} and the full training history is
intact.}), the existing \emph{analysis workflow} (generating
\texttt{peak\_alloc\_events\_GPU0.json}) did not perform a full trace replay.
Instead it relied on \texttt{snapshot.segments}---the point-in-time
current-state view---which only shows 155 blocks at 941.67\,MB, still
short of the 1173.02\,MB true peak.

%─────────────────────────────────────────────────────────────────────────────
\subsection{Logical Error 2: \texttt{peak\_alloc\_events} Is a Subset,
            Not the Total Peak}
\label{sec:err2subset}
%─────────────────────────────────────────────────────────────────────────────

The \texttt{process\_alloc\_data} function in
\texttt{ref\_find\_max/MemoryViz.js} classifies every allocation event
into one of two categories:

\begin{enumerate}
  \item \textbf{\texttt{draw\_elem}}: the \texttt{max\_entries} \emph{largest}
    allocations, tracked individually as plotted polygons.
  \item \textbf{\texttt{summarized\_mem}}: all remaining (smaller) allocations,
    aggregated into a single running counter \texttt{total\_summarized\_mem}.
\end{enumerate}

The \textbf{peak detection} correctly sums both:
\begin{lstlisting}[language=JavaScript, caption={Correct total-memory peak
  check inside the action loop (\texttt{MemoryViz.js}).}]
const current_total = total_mem + total_summarized_mem;  // includes BOTH
if (current_total > max_size) {
    max_size  = current_total;   // TRUE peak: large + summarized
    peak_elem = elem;
}
\end{lstlisting}

However, the \textbf{output list} \texttt{peak\_alloc\_events} is built by
filtering \texttt{data}---which only contains \texttt{draw\_elem} entries---at
the peak timestep:

\begin{lstlisting}[language=JavaScript, caption={Incomplete peak-events
  collection: summarized blocks are silently excluded.}]
const peak_allocs = data.filter(d => {
    if (d.elem === 'summarized') return false;  // <-- excludes summarized!
    const ts = d.timesteps;
    return ts[0] <= peak_timestep && peak_timestep <= ts[ts.length - 1];
});
const peak_alloc_events = peak_allocs.map(d => elements[d.elem]);
// peak_alloc_events contains ONLY draw_elem blocks -- NOT the total peak
\end{lstlisting}

Consequently, when \texttt{peak\_alloc\_events\_GPU0.json} is generated from
this output and its block sizes are summed, the result (625.77\,MB) is the
contribution of the \emph{individually tracked large blocks only}.
The many small blocks in \texttt{summarized\_mem}---LayerNorm parameter
tensors, bias vectors, scalar loss tensors, pipeline communication buffers,
and similar small allocations---are \textbf{not included}.
Their combined contribution at the true peak timestep accounts for the
remaining:
\begin{equation}
  \underbrace{1173.02\,\text{MB}}_{\text{true total peak}}
  - \underbrace{625.77\,\text{MB}}_{\text{draw\_elem only (JSON)}}
  \approx 547\,\text{MB of omitted summarized blocks}
  \label{eq:missing}
\end{equation}

\begin{bugbox}
\textbf{Summary of Error 2.}
\texttt{peak\_alloc\_events\_GPU0.json} contains only the
\texttt{draw\_elem} (large-block) subset of the peak, not the
complete memory picture.  Treating this JSON as the total peak memory is
incorrect; the correct total ($\mathtt{max\_size}$) is 1173.02\,MB.
\end{bugbox}

%─────────────────────────────────────────────────────────────────────────────
\subsection{Where Does the True Maximum Occur?}
\label{sec:wheremax}
%─────────────────────────────────────────────────────────────────────────────

\paragraph{Training-step context.}
The true peak (event 8015 in \texttt{device\_traces[0]}) occurs
\emph{inside} a training step during the \textbf{backward pass with
GPipe activation recomputation}.  The memory timeline exhibits the following
salient feature: immediately after event 8015, a 147.2\,MB block is freed
at event 8091 (a drop of $-154{,}337{,}280$ bytes), reducing total memory to
$941.67\,\text{MB}$.  This 147.2\,MB block corresponds to an
embedding-sized tensor ($\approx V \times d \times 4\,\text{B}
= 50{,}257 \times 768 \times 4 = 154\,\text{MB}$) that GPipe's
\texttt{checkpoint='always'} recomputes as part of the
\texttt{EmbeddingBlock} forward pass during backward, holds in memory
while gradients propagate, then frees.

\paragraph{Components alive at the true peak (event 8015).}

\begin{table}[h]
\centering
\caption{Memory components simultaneously alive at event 8015
(GPU\,0, second training step).}
\label{tab:peak_components}
\begin{tabular}{lrl}
\toprule
\textbf{Component} & \textbf{Approx.\ size} & \textbf{Notes} \\
\midrule
Model parameters ($\theta$)            & $\sim214\,\text{MB}$
  & Persistent; large vocab embedding \\
AdamW \texttt{exp\_avg} ($m_t$)        & $\sim214\,\text{MB}$
  & Persistent; allocated in step 1 \\
AdamW \texttt{exp\_avg\_sq} ($v_t$)    & $\sim214\,\text{MB}$
  & Persistent; allocated in step 1 \\
Gradients ($g_t$)                      & $\sim214\,\text{MB}$
  & Re-created each step by \texttt{backward()} \\
\midrule
\textbf{Persistent subtotal}           & $\mathbf{\sim856\,\text{MB}}$
  & Matches observed final state (941\,MB) \\
\midrule
GPipe recomputed embedding activation  & $\sim147\,\text{MB}$
  & Freed at event 8091; \textbf{not in final snapshot} \\
Other transient backward intermediates & $\sim{170}\,\text{MB}$
  & Freed during step; \textbf{not in final snapshot} \\
\midrule
\textbf{True peak total}               & $\mathbf{1173.02\,\text{MB}}$
  & From trace replay of \texttt{device\_traces[0]} \\
\bottomrule
\end{tabular}
\end{table}

\noindent The true peak \emph{cannot be seen} by inspecting
\texttt{snapshot.segments} alone---all transient components above the dashed
line in Table~\ref{tab:peak_components} have already been freed by the time
\texttt{\_dump\_snapshot} is called.  A correct peak analysis requires
\textbf{full trace replay of} \texttt{device\_traces}, as implemented in
\texttt{MemoryViz.js}'s \texttt{process\_alloc\_data}, but crucially the
caller must read \texttt{max\_size} (the algorithm's total peak including
summarized blocks), \emph{not} the incomplete \texttt{peak\_alloc\_events}
list.

%─────────────────────────────────────────────────────────────────────────────
\subsection{Corrected Peak Memory Formula for GPU\,0}
\label{sec:correctedformula}
%─────────────────────────────────────────────────────────────────────────────

Including the transient GPipe activation recomputation tensors, the correct
peak for GPU\,0 is:
\begin{equation}
  M_{\text{true peak}}^{\text{GPU0}}
  = \underbrace{P_0}_{\theta}
  + \underbrace{P_0}_{g_t}
  + \underbrace{P_0 + P_0}_{m_t + v_t}
  + \underbrace{P_0}_{\sqrt{v_t}\text{ (AdamW temp.)}}
  + \underbrace{A_{\text{recompute}}}_{\text{GPipe activation}}
  \label{eq:truepeak_formula}
\end{equation}
where $P_0 \approx 214\,\text{MB}$ (GPU\,0 parameters, dominated by the
$50{,}257 \times 768$ vocabulary embedding) and
$A_{\text{recompute}} \approx 317\,\text{MB}$ (transient activation tensors
during backward recomputation, including the 147\,MB embedding-sized block
and $\sim 170\,\text{MB}$ of smaller intermediates).

\begin{equation}
\boxed{
  M_{\text{true peak}}^{\text{GPU0}}
  = 5 \times 214 + 317
  \approx 1387\,\text{MB}
  \quad (\text{theoretical})
}
\end{equation}

\noindent The theoretical estimate ($\approx 1387\,\text{MB}$) is close to
but somewhat higher than the measured trace value (1173\,MB) because the
AdamW \texttt{sqrt\_buf} and the recomputed activation tensors do not all
reach their peak \emph{simultaneously}: the snapshot captures the global
maximum of the running total, which peaks slightly below the sum of all
theoretical components.

%─────────────────────────────────────────────────────────────────────────────
\subsection{Summary: Why 282 Blocks / 625.77\,MB Is Wrong}
\label{sec:wrongsummary}
%─────────────────────────────────────────────────────────────────────────────

\begin{table}[h]
\centering
\caption{Root causes of the 282-block underestimate.}
\label{tab:errors}
\begin{tabular}{clp{8cm}}
\toprule
\textbf{\#} & \textbf{Error} & \textbf{Consequence} \\
\midrule
1 & Snapshot taken after training loop
  & \texttt{snapshot.segments} shows 155-block final state (941\,MB),
    not the 193-block peak (1173\,MB).  All transient backward tensors
    are already freed. \\[4pt]
2 & \texttt{peak\_alloc\_events} excludes \texttt{summarized\_mem}
  & The 282-block JSON lists only \texttt{draw\_elem} (large) blocks
    (625.77\,MB); the $\sim 547\,\text{MB}$ of small summarized blocks
    at the same peak timestep are silently omitted. \\[4pt]
3 & Per-GPU loop in \texttt{start\_memory\_history\_all\_gpus()}
  & \texttt{\_record\_memory\_history()} is a global API; the loop
    calls it 4 times unnecessarily.  This is harmless to the ring buffer
    here (no overflow), but demonstrates the API misunderstanding that
    contributed to Error 1. \\
\bottomrule
\end{tabular}
\end{table}

\noindent The combined effect of Errors 1 and 2 is that the analysis
attributes a peak of $625.77\,\text{MB}$ to GPU\,0, while the true maximum
obtained from a correct trace replay is $\mathbf{1173.02\,\text{MB}}$---
nearly \textbf{double} the reported figure.  Because the true peak occurs
during the backward pass with GPipe activation recomputation (not during
\texttt{optimizer.step()}) the AdamW 3X story from
Section~\ref{sec:3x} is only \emph{part} of the explanation: the dominant
factor at the system level is the recomputed embedding activation of
$\approx 147\,\text{MB}$ that coexists with the full $5P$
optimizer-step peak.

%=============================================================================
\section{Conclusion}
\label{sec:conclusion}
%=============================================================================

\begin{enumerate}
  \item AdamW is instantiated in \textbf{two places}: the isolated
        test script \texttt{why\_adam\_times\_3\_here/iiiandadam.py} and
        the full GPT-2 profiler \texttt{gpt2\_gpipe\_memory\_profile.py},
        both using \texttt{torch.optim.AdamW}.

  \item The two \emph{persistent} optimizer states (\texttt{exp\_avg},
        \texttt{exp\_avg\_sq}) cost \textbf{2X} of the model size after
        \texttt{step()} completes.

  \item The Adam update formula requires $\sqrt{v_t}$, which is computed
        as a \emph{new tensor} (not in-place).  This temporary buffer adds
        another \textbf{1X} during the step, bringing the \textbf{peak
        AdamW overhead to 3X}.

  \item Because this analysis focuses on \textbf{maximum (peak) memory},
        the correct figure is \textbf{3X}.  Post-step persistence,
        \texttt{zero\_grad()}, and whether a tensor is ``kept'' after the
        step are all irrelevant to the peak.

  \item The \textbf{peak-detection algorithm} in
        \texttt{ref\_find\_max/MemoryViz.js} replays the alloc/free event
        trace and identifies the timestep of maximum total allocation,
        correctly capturing transient tensors.

  \item A \textbf{code bug} exists in \texttt{gpt2\_gpipe\_memory\_profile.py}:
        the per-GPU loop in \texttt{start\_memory\_history\_all\_gpus()} and
        \texttt{stop\_memory\_history\_all\_gpus()} is unnecessary because
        \texttt{torch.cuda.memory.\_record\_memory\_history()} is a single
        global call in modern PyTorch.
\end{enumerate}

\end{document}
